\chapter{Introducción al CAD y Entorno de Trabajo}

\section{Fundamentos del CAD y Entorno de Trabajo}

El Dibujo Técnico Asistido por Computadora (CAD, por sus siglas en
inglés: Computer-Aided Design) es una tecnología que permite la
creación, modificación, análisis y optimización de diseños utilizando
computadoras. En el nivel secundario, es la herramienta fundamental
para el dibujo técnico, reemplazando tableros y reglas por entornos
virtuales precisos.

El CAD ha revolucionado la ingeniería y la arquitectura. Mientras que
el dibujo técnico manual depende de la destreza física del dibujante,
el CAD garantiza precisión, repetibilidad y facilidad de modificación.

\subsection{Historia y evolución}

El origen del CAD se remonta a la década de 1960 con el sistema
Sketchpad de Ivan Sutherland, que permitió por primera vez la
interacción gráfica con una computadora. A lo largo de las décadas, el
software CAD evolucionó de herramientas de dibujo 2D (líneas y arcos)
a potentes sistemas de modelado 3D paramétrico, que permiten simular
el comportamiento físico de las piezas.

\subsection{El entorno de trabajo CAD}

El entorno de trabajo se basa en un plano cartesiano infinito donde
cada punto se define por coordenadas $(x, y, z)$. La precisión es
absoluta, permitiendo dibujar unidades de medida exactas (milímetros,
metros) que se ajustan mediante escalas de impresión.

Un software CAD típico consta de los siguientes elementos:

\begin{itemize}
	\item \textbf{Barra de herramientas (Ribbon):} Acceso rápido a los
	      comandos más usados (línea, círculo, recortar).
	\item \textbf{Línea de comandos (Command Line):} El corazón del
	      sistema, donde se introducen órdenes mediante el teclado y se
	      reciben mensajes del software.
	\item \textbf{Espacio Modelo:} Espacio infinito donde se dibuja en
	      escala real (1:1).
	\item \textbf{Espacio Papel (Layout):} Entorno de diseño para
	      configurar la lámina de impresión con sus vistas, escala y
	      rótulos.
	\item \textbf{Capas (Layers):} Herramienta que permite organizar el
	      dibujo en distintos niveles, gestionando colores, grosores de
	      línea y visibilidad.
\end{itemize}

\subsection{Tipos de software CAD}

Existen dos tipos principales de CAD:

\begin{enumerate}
	\item \textbf{CAD 2D (Drafting):} Centrado en la representación
	      bidimensional, es el sucesor digital del dibujo en tablero.
	\item \textbf{CAD 3D (Modelado):} Permite la creación de sólidos y
	      superficies, fundamentales para ingeniería mecánica y diseño
	      industrial. Se divide a su vez en:
	      \begin{itemize}
		      \item \textbf{Paramétrico:} Los cambios en una dimensión
		            actualizan automáticamente todo el modelo.
		      \item \textbf{Directo:} Se manipulan las geometrías directamente
		            sin historial de operaciones.
	      \end{itemize}
\end{enumerate}

\section{Cuestionario}
\begin{enumerate}
	\item ¿Qué significan las siglas CAD?
	\item ¿Por qué es importante el uso de coordenadas cartesianas en CAD?
	\item ¿Cuál es la diferencia entre el Espacio Modelo y el Espacio Papel?
	\item Define qué es una capa (layer) y para qué sirve.
	\item ¿Cómo afecta la escala 1:1 en el diseño de un objeto pequeño?
	\item ¿Por qué la precisión digital supera al dibujo manual?
	\item ¿Quién creó el primer sistema CAD y en qué década?
	\item Explica la diferencia entre CAD paramétrico y directo.
\end{enumerate}

\section{Parte Práctica}
\begin{enumerate}
	\item Configurar una unidad de dibujo en milímetros en el software.
	\item Crear una línea de 100mm horizontal en el eje X.
	\item Crear un cuadrado de 50x50mm usando coordenadas relativas.
	\item Dibujar un círculo de diámetro 40mm centrado en el origen.
	\item Crear dos capas nuevas: "Ejes" (roja, línea discontinua) y "Contornos" (negra, línea continua).
	\item Mover el círculo a la capa "Contornos".
	\item Dibujar un rectángulo de 100x60mm y centrarlo en el origen.
	\item Aplicar un zoom a extensión para ver todo el dibujo.
	\item Crear un nuevo archivo y configurar un tamaño de papel A4.
	\item Realizar un dibujo libre usando al menos 3 capas diferentes.
\end{enumerate}
